% BHU PYQ -- Transcribed & Corrected
% Paper: HINMD21: साहित्य मिथक एवं इतिहास
% Exam: Under Graduate Semester - II Examination (NEP) 2024-25
% Ref Code: CE/CONF/D-II/124
% Category: Multidisciplinary Course (MDC)
% Department: Hindi / Faculty of Arts

\documentclass[11pt,a4paper]{article}
\usepackage[a4paper,top=2.5cm,bottom=2.5cm,left=2.8cm,right=2.8cm,headheight=14pt]{geometry}
\usepackage{amsmath,amssymb}
\usepackage{fontspec}
\setmainfont{Kohinoor Devanagari}
\usepackage{microtype,fancyhdr,enumitem,hyperref,lastpage,parskip}

\hypersetup{
    colorlinks=true,
    urlcolor=black,
    linkcolor=black,
    pdftitle={HINMD21: Sahitya Mithak Evam Itihas -- BHU Sem II 2024-25 (NEP MDC)}
}

\pagestyle{fancy}
\fancyhf{}
\renewcommand{\headrulewidth}{0.4pt}
\renewcommand{\footrulewidth}{0.4pt}
\fancyhead[L]{\small \textit{Banaras Hindu University}}
\fancyhead[R]{\small \textit{HINMD21: साहित्य मिथक एवं इतिहास (2024-25)}}
\fancyfoot[C]{\small Page \thepage\ of \pageref{LastPage}}

\newcommand{\pts}[1]{\hfill \textit{[#1]}}

\begin{document}

\begin{center}
  {\large\bfseries BANARAS HINDU UNIVERSITY}\\[2pt]
  {\normalsize Under Graduate Semester II Examination (NEP) 2024-25}\\[6pt]
  \rule{0.8\linewidth}{0.5pt}\\[6pt]
  {\large\bfseries HINDI (MDC)}\\[4pt]
  {\large\bfseries Paper Code: HINMD-21 : साहित्य मिथक एवं इतिहास}\\[2pt]
  {\small Ref. No.: CE/CONF/D-II/124}\\[4pt]
  {\normalsize Time: 2$\frac{1}{2}$ Hours\hfill Maximum Marks: 70}
\end{center}

\rule{\linewidth}{0.4pt}

\smallskip

\noindent \textit{अनुदेश : निर्देशानुसार प्रश्नों के उत्तर दीजिए। प्रश्नों के अंक उनके सम्मुख अंकित हैं।}

\bigskip

\noindent \textbf{प्रश्न 1.} निम्नलिखित दीर्घउत्तरीय प्रश्नों में से किन्हीं तीन प्रश्नों के उत्तर दीजिए: \pts{3 $\times$ 10 = 30}
\begin{enumerate}[label=(\alph*)]
  \item 'झाँसी की रानी' उपन्यास के आधार पर रानी लक्ष्मीबाई का चरित्र-चित्रण कीजिए।
  \item दया प्रकाश सिन्हा के नाटक 'सम्राट अशोक' की कथावस्तु में मौजूद कल्पना और यथार्थ तत्व को रेखांकित कीजिए।
  \item 'महाराणा का महत्त्व' शीर्षक गीति-नाट्य के कथ्य और भाषा पर एक निबंध लिखिए।
  \item 'कैकेयी अनुताप' काव्य के आधार पर भारतीय परिवार संबंधों और मूल्यों के सांस्कृतिक महत्त्व पर अपने विचार लिखिए।
  \item 'महाराणा का महत्त्व' गीति-नाट्य में राणा प्रताप के चरित्र की कौन-सी विशेषताएँ वर्णित हैं, जिनका भारतीय संस्कृति और राष्ट्रीय आंदोलन की दृष्टि से विशेष महत्त्व है?
\end{enumerate}

\bigskip

\noindent \textbf{प्रश्न 2.} निम्नलिखित विषयों में से किन्हीं चार पर टिप्पणी लिखिए: \pts{4 $\times$ 4 = 16}
\begin{enumerate}[label=(\alph*)]
  \item 'सम्राट अशोक' नाटक की रंगमंचीयता
  \item 'झाँसी की रानी' उपन्यास में मौजूद राष्ट्रीय-भावना
  \item मिथक और इतिहास में अंतर्संबंध
  \item 'महाराणा का महत्त्व' का मूल प्रतिपाद्य
  \item 'कैकेयी अनुताप' का सभा-दृश्य
  \item नाटक के तत्व
  \item गीतिनाट्य और जयशंकर प्रसाद
  \item हिन्दी नवजागरण और राष्ट्रवाद
\end{enumerate}

\bigskip

\noindent \textbf{प्रश्न 3.} निम्नलिखित अवतरणों में से किन्हीं चार अंशों की सप्रसंग व्याख्या कीजिए: \pts{6 $\times$ 4 = 24}
\begin{enumerate}[label=(\alph*)]
  \item अब इस देश का भाग्य लौट गया है। अंग्रेजों के भाग्य का सूर्योदय हुआ है। उन लोगों के प्रताप के सामने यहाँ के सब जन निस्तेज हो गये हैं।
  \item मैं डरपोक कभी नहीं हो सकती। आप कहते हैं-मन्नू तू ताराबाई बनना, जीजाबाई और सीता होना। यह सब भुलावा क्यों? अथवा क्या ये सब डरपोक थीं? बेटी, ये सब सती और वीर थीं, परन्तु समय बदलता रहता है। बदल गया है।
  \item कानन में पतझड़ भी कैसा फैल के भीषण निज आतंक दिखाता था,\\
  कड़े सूखे पत्तों के ही 'खड़-खड़' शब्द से अपना कुत्सित क्रोध प्रकट था कर रहा\\
  हरे-हरे द्रुम दल को खूब लथेड़ता घूम रहा था, क्रूर सदृश उस भूमि में।
  \item हम बंदी हैं। (अपने ऊपर हसते हुए) राजा अशोक बंदी है। (चारों ओर देखकर) तो यह बंदीगृह है। (नंदिता से) तुम तो बंदी नहीं हो। तुम इस बंदीगृह में क्या कर रही हो। (चिल्लाते हुए) निकलो यहाँ से। (नंदिता घबराकर भागती हुई बाहर निकल जाती है। अशोक क्रोध में मुट्ठियाँ तानते हुए शेर की तरह इधर-उधर चक्कर लगाता है। महामंत्री का प्रवेश।)
  \item देव दिवाकर भी असहय थे हो रहे यह छोटा-सा झुंड सहन कर ताप को,\\
  बढ़ता ही जाता है अपने मार्ग में। शिविका को घेरे थे वे सैनिक सभी\\
  जो गिनती में शत थे, प्रण में वीर थे।
  \item यह सच है तो अब लौट चलो तुम घर को। चौंके सब सुनकर अटल केकयी-स्वर को।\\
  सबने रानी की ओर अचानक देखा, वैधव्य-तुषारावृता यथा विधु-लेखा।
  \item यदि मैं उकसाई गई भरत से होऊँ, तो पति समान ही स्वयं पुत्र भी खोऊँ।\\
  ठहरो, मत रोको मुझे, कहूँ सो सुन लो। पाओ यदि उसमें सार उसे सब चुन लो।\\
  करके पहाड़-सा पाप मौन रह जाऊँ? राई भर भी अनुताप न करने पाऊँ?
\end{enumerate}

\bigskip
\begin{center}
  \textbf{***}
\end{center}

\end{document}
